\chapter{核心模块设计与实现}

本章详细阐述系统四个核心模块——掩模生成器、噪声注入器、渲染驱动器、时序控制器——以及辅助模块的设计思路与关键实现。系统以 Python 实现，使用 NumPy 进行数值计算、PyCryptodome 提供 ChaCha20、moderngl 提供 GPU 着色器渲染，并配套 118 项单元测试保证算法正确性。

\section{掩模生成器}

掩模生成器负责由密钥 $K$ 和周期序号生成随机点阵掩模与子帧置换序列，是系统安全性的核心。其实现包含密钥派生、无偏索引生成、均匀性检验三个关键环节。

\subsection{周期子密钥派生}

为保证相邻周期掩模统计独立，模块以主密钥 $K$ 和周期序号 $t$ 派生周期子密钥。实现上采用 HMAC-SHA256 构造 HKDF 风格的派生：以主密钥为密钥、周期序号与标签拼接为消息计算 HMAC，取摘要前 32 字节作为该周期的 ChaCha20 子密钥。这样即使攻击者获取了某周期掩模，也无法推断其他周期，对应技术方案中"密钥派生多样性"的抗自适应攻击设计。

\subsection{无偏随机索引生成}

索引矩阵 $R(x,y) \in \{0,1,\dots,n-1\}$ 由 ChaCha20 密钥流生成。一个容易被忽视的问题是取模偏置：若直接对均匀字节做 \verb|% n|，当 $n$ 不是 2 的幂时会引入系统性偏差（例如 $256 \bmod 3 = 1$，使值 0 多出现一次）。本系统采用\textbf{拒绝采样}消除该偏置：读取 $\lceil\log_2 n\rceil$ 位整数，若落在 $[n, 2^{\lceil\log_2 n\rceil})$ 区间则丢弃重采，从而保证严格均匀分布。核心逻辑如代码所示。

\begin{quote}\small\ttfamily
bits = ceil(log2(n));\ ceil\_pow2 = 1 << bits \\
vals = chacha20\_stream() \& (ceil\_pow2 - 1) \\
accepted = vals[vals < n]\quad\# 拒绝 [n, 2\^{}bits) 区间
\end{quote}

\subsection{均匀性检验}

为对应技术方案"步骤D（均匀性验证）"，模块对生成的索引矩阵执行卡方检验：统计每个取值 $k$ 的像素数量，与理论期望 $W\cdot H/n$ 比较，若 $p > 0.01$ 则认为分布均匀，否则偏移种子重新生成（最多重试 5 次）。生成的 $n$ 个布尔掩模严格满足完备性 $\sum_k M_k = \mathbf{1}$，该约束由单元测试逐像素断言。此外，模块支持预生成未来若干周期的掩模与置换序列并存入环形缓冲，对应技术方案中"预生成 1024 个周期"的设计。

\section{噪声注入器}

噪声注入器负责生成时域互补的对抗子噪声。模块支持多种梯度来源：针对 EasyOCR 识别器的可微端到端梯度、轻量影子模型梯度、OCR 对比度启发式梯度，当真实梯度不可用时自动回退到代理梯度并记录来源，保证实验可复现。

基础噪声生成后，通过 \verb|split_complementary| 将其分解为满足 $\sum_k N_k = 0$ 的 $n$ 个互补子噪声，分解残差由单元测试验证为精确 0。模块还实现了 \verb|split_complementary_spatial|，在时域互补的基础上额外约束相邻像素的棋盘式极性抵消，对应技术方案"空间-时间联合扰动"。为对抗 AI 自适应攻击，模块维护针对 Tesseract、EasyOCR、PaddleOCR、YOLOv8、Faster R-CNN 五种目标模型的噪声模板，运行时以伪随机序列轮换选择，并支持在线监测识别率、动态调整扰动预算与频域分布。

\section{渲染驱动器}

渲染驱动器在像素层面将掩模、子噪声、亮度补偿合成为子帧。模块对每个子帧执行：读取掩模 $M_k$ 与子噪声 $N_k$，计算 $I_k^{\text{raw}} = I \cdot M_k + N_k$，应用亮度补偿，必要时在周期末生成反色帧，最终输出。

考虑到部署环境的图形栈差异，模块采用 GPU 优先、软件回退的双路径设计：工厂函数 \verb|create_renderer| 优先初始化 moderngl 的 GPU 着色器渲染器，若 moderngl 未安装或 OpenGL 上下文创建失败（如无头环境、远程桌面），则捕获异常并自动回退到等价的 NumPy 软件渲染器，二者输出经单元测试验证为像素级一致。这一设计保证了原型在各种环境下都能运行，但也意味着\textbf{实际渲染路径可能是 CPU 软件渲染}——这一点在性能评测（第五章）中如实反映，本原型不直接承诺技术方案中假设 GPU Compute Shader 并行的低开销指标。

\begin{quote}\small\ttfamily
renderer = GPURenderer(...) \\
try:\ renderer.\_init\_gl();\ return renderer \\
except\ ImportError:\ \# moderngl 缺失 \\
except\ Exception:\ \# GL 上下文失败 \\
\quad return SoftwareRenderer(...)\quad\# 等价回退
\end{quote}

\section{时序控制器}

时序控制器负责子帧的时序调度与同步。模块实现了三项关键机制。其一，伪随机置换：每周期以 Fisher-Yates 算法对子帧输出顺序进行无偏洗牌，使输出序列不固定为 $1\to2\to\dots\to n$，提高攻击者完整周期对齐的难度。其二，时序令牌：维护包含周期编号、置换序列哈希、VBlank 计数的时序令牌，以锁保护原子更新，对应技术方案中防止多线程竞争的"时序令牌"设计。其三，应急黑帧：当检测到渲染未在 VBlank 前完成时，输出黑帧而非未完成的中间状态，防止半成品画面被相机捕获。

此外，模块实现了显示接口带宽计算 $\text{Bandwidth} = H\cdot V\cdot f_r\cdot \text{Bpp}\cdot 1.2$ 及对 DP1.4/DP2.0/HDMI2.1 容量的判断。需要说明，本原型的时序为\textbf{软件模拟}，真实的硬件 VBlank 中断、实时线程调度与 $\pm0.5$ms 级硬同步需操作系统驱动层支持，属未来工作。

\section{辅助模块}

\textbf{HDR 亮度补偿}模块实现了 PQ（ST.2084）与 HLG 编解码、线性 RGB 与 ICtCp 色彩空间互转、超色域软裁剪以及 $L_k\cdot n$ 的 HDR 亮度提升，并以环境光自适应闭环根据模拟光照调整背光与补偿系数，对应技术方案 4.3 节。

\textbf{多显示器同步}模块实现主从 VBlank 广播与异构刷新率的最小公倍数虚拟时钟调度，单元测试验证同步误差最大 0.05ms，并校验任意相机曝光窗口内无完整帧泄露。

\textbf{视觉疲劳策略}模块实现了静/动态内容自适应刷新率、蓝光抑制色温调整、观看距离亮度补偿三项纯函数策略，对应技术方案 7.4 节。

\textbf{配置持久化}模块以数据类封装拆分因子、噪声预算、反色系数、密钥等参数，支持 JSON 读写并接入主程序，对应技术方案 5.1 节的用户配置加载。

\section{本章小结}

本章详细阐述了四个核心模块的设计与实现：掩模生成器通过 ChaCha20、HKDF 派生与拒绝采样实现无偏、周期独立、分布均匀的随机点阵掩模；噪声注入器实现了多梯度源对抗噪声与时域互补分解；渲染驱动器采用 GPU 优先、软件回退的双路径合成；时序控制器实现了伪随机置换、时序令牌与应急黑帧。同时介绍了 HDR 补偿、多显示器同步等辅助模块。本章如实标注了软件模拟时序、CPU 渲染回退等与产品级实现的差异，为第五章基于真实实现的实验评测提供了准确基础。
